%%%%%%%%%%%%%%%%%%%%%%%%%%%%%%%%%%%%%%%%%%%%%%%%%%%%%%%%%%%%%%%%
%  chapter1.tex  --  Introduction
%  ISOLA Game in Prolog with Alpha-Beta Algorithm
%  Artificial Intelligence, ISEL 2025/2026
%%%%%%%%%%%%%%%%%%%%%%%%%%%%%%%%%%%%%%%%%%%%%%%%%%%%%%%%%%%%%%%%

\chapter{Introduction}
\label{chap:intro}

\section{The ISOLA Game}
\label{sec:game}

\textit{ISOLA} (also known as \textit{Isolation} or \textit{Solitude}) is a
two-player strategy game created by Bernd Kienitz in 1972~\cite{kienitz1972isola}. The objective is
straightforward: isolate the opponent's pawn such that it cannot make any
valid move on its turn.

\subsection{Game Rules}
\label{subsec:rules}

The board is a square grid (typically 6{\texttimes}6). Each player has a pawn that
starts at opposite corners:

\begin{itemize}
  \item \textbf{Player~1} (\texttt{x}) starts at position
        (row~0, column~0) --- top-left corner.
  \item \textbf{Player~2} (\texttt{o}) starts at position
        (row~5, column~5) --- bottom-right corner.
\end{itemize}

On each turn, a player must perform exactly \textbf{two actions}:

\begin{enumerate}
  \item \textbf{Move} their pawn to an adjacent free square
        (horizontally, vertically, or diagonally --- like a king in
        chess, but only one step at a time).
  \item \textbf{Remove} any empty square from the board (except the
        squares occupied by the pawns), marking it as permanently
        inaccessible (\texttt{\#}).
\end{enumerate}

\textbf{It is not permitted} to jump over the opponent or to enter
already-removed squares. The player who leaves the opponent with no
valid move at the start of their turn wins.

\subsection{Game Example}
\label{subsec:example}

The example below illustrates the progression of a game on a 6{\texttimes}6 grid.
In the initial state, Player~1 moves from (0,0) to (1,1) and removes
square (2,2).

\begin{figure}[htbp]
  \centering
  \begin{minipage}[t]{0.45\textwidth}
    \centering
    \textbf{Initial state}
    \begin{verbatim}
       0  1  2  3  4  5
    +------------------
 0  |  x  .  .  .  .  .
 1  |  .  .  .  .  .  .
 2  |  .  .  .  .  .  .
 3  |  .  .  .  .  .  .
 4  |  .  .  .  .  .  .
 5  |  .  .  .  .  .  o
    \end{verbatim}
  \end{minipage}
  \hfill
  \begin{minipage}[t]{0.45\textwidth}
    \centering
    \textbf{After P1 moves (0,0)$\to$(1,1), removes (2,2)}
    \begin{verbatim}
       0  1  2  3  4  5
    +------------------
 0  |  .  .  .  .  .  .
 1  |  .  x  .  .  .  .
 2  |  .  .  #  .  .  .
 3  |  .  .  .  .  .  .
 4  |  .  .  .  .  .  .
 5  |  .  .  .  .  .  o
    \end{verbatim}
  \end{minipage}
  \caption{Progression of an ISOLA 6{\texttimes}6 game.
           The symbol \texttt{\#} represents a permanently removed square.}
  \label{fig:game_example}
\end{figure}

The strategy consists of keeping one's own pawn with many movement options
(\textbf{mobility}) while progressively reducing the opponent's mobility,
ultimately isolating them completely.

% -------------------------------------------------------------------

\section{Learning Objectives}
\label{sec:objectives}

This work aims to consolidate the following learning objectives:

\begin{enumerate}
  \item Understand and master the syntax of Prolog programs.
  \item Use basic constructs: atoms, compound terms, lists,
        operators, and arithmetic.
  \item Understand and control \textit{backtracking}.
  \item Be able to use typical built-in predicates
        (\texttt{findall/3}, \texttt{nth0/3}, \texttt{maplist/3}, etc.).
  \item Implement a two-player game with perfect information.
  \item Understand the \textit{minimax} principle.
  \item Understand the \textit{alpha-beta} algorithm: an efficient
        implementation of \textit{minimax}.
\end{enumerate}

% -------------------------------------------------------------------

\section{Introduction to Prolog}
\label{sec:prolog_intro}

To understand the implementation, it is essential to understand the
fundamentals of Prolog~\cite{bratko2001prolog, sterling1994art}. Unlike languages such as Java or Python, Prolog does
not describe \emph{how} to compute a result, but rather \emph{what is true}.
The runtime system is responsible for determining how to derive the answer.

\subsection{Analogy: The Detective and the Clues}

Consider Prolog as a detective. One provides it with a set of
\textbf{facts} (known clues) and \textbf{rules} (logical deductions),
and then poses \textbf{queries} (questions). The detective systematically
searches all possible combinations until an answer is found.

\begin{lstlisting}[language=Prolog,
    caption={Basic facts and rules in Prolog},
    label={lst:prolog_basics},
    numbers=left,
    basicstyle=\small\ttfamily,
    frame=single]
% Facts: who likes what
likes(mary, cats).
likes(john, dogs).
likes(mary, dogs).

% Rule: two people have something in common
in_common(X, Y, Thing) :-
    likes(X, Thing),
    likes(Y, Thing),
    X \= Y.

% Query: ?- in_common(mary, john, Thing).
% Answer: Thing = dogs.
\end{lstlisting}

The clause beginning with \texttt{:-} is read as ``if'': \texttt{in\_common(X,\,Y,\,Thing)}
is true \textbf{if} \texttt{likes(X,\,Thing)} is true \textbf{and}
\texttt{likes(Y,\,Thing)} is true \textbf{and} \texttt{X} is not equal to \texttt{Y}.

\subsection{Terms and Unification}

The central element of Prolog is the \textbf{term}. There are four types:

\begin{itemize}
  \item \textbf{Atom} --- symbolic constant: \texttt{cat}, \texttt{x},
        \texttt{'.'}, \texttt{true}.
  \item \textbf{Number} --- integer or floating-point: \texttt{3}, \texttt{4.5}.
  \item \textbf{Variable} --- begins with an uppercase letter or underscore:
        \texttt{X}, \texttt{Player}, \texttt{\_}.
  \item \textbf{Compound term} --- functor with arguments:
        \texttt{move(1/2,\,3/4)}, \texttt{[a,b,c]}.
\end{itemize}

\textbf{Unification} is the mechanism by which Prolog binds values to
variables. When one writes \texttt{X~=~3}, Prolog unifies \texttt{X} with
\texttt{3}. Unification succeeds only if the two terms share the same
structure.

\begin{lstlisting}[language=Prolog,
    caption={Unification examples},
    label={lst:unification},
    numbers=left,
    basicstyle=\small\ttfamily,
    frame=single]
?- X = 5.            % X = 5  (succeeds)
?- 2/3 = R/C.        % R = 2, C = 3 (term decomposition)
?- [H|T] = [1,2,3].  % H = 1, T = [2,3] (head and tail)
?- 2 = 3.            % false (unification fails)
\end{lstlisting}

\subsection{Backtracking}
\label{subsec:backtracking}

\textit{Backtracking} is Prolog's search mechanism.
When an attempt to satisfy a goal fails, Prolog
\textbf{backtracks} automatically and tries the next available
alternative.

\textbf{Analogy}: imagine a maze. Prolog tries one path;
if it reaches a dead end, it backtracks and tries the next
corridor, until all paths have been explored.

\begin{lstlisting}[language=Prolog,
    caption={Backtracking demonstration},
    label={lst:backtrack},
    numbers=left,
    basicstyle=\small\ttfamily,
    frame=single]
color(red).
color(green).
color(blue).

% Query: ?- color(X).
% X = red    (first solution)
% X = green  (after requesting more: ";")
% X = blue   (third solution)
% false      (no more solutions)
\end{lstlisting}

The predicate \texttt{findall/3} exploits \textit{backtracking} to
collect \emph{all} solutions into a list, which is used extensively
throughout this project:

\begin{lstlisting}[language=Prolog,
    caption={findall collects all solutions},
    label={lst:findall_intro},
    numbers=left,
    basicstyle=\small\ttfamily,
    frame=single]
?- findall(X, color(X), List).
% List = [red, green, blue]
\end{lstlisting}

\subsection{Lists in Prolog}

In Prolog, a list is represented as
\texttt{[Elem1,\,Elem2,\,\ldots]}. The notation \texttt{[H|T]} separates
the head (\texttt{H}) from the tail (\texttt{T}).

In the ISOLA game, the board is a \textbf{list of lists}: each
row of the board is a list of atoms.

\begin{lstlisting}[language=Prolog,
    caption={3{\texttimes}3 board as a list of lists},
    label={lst:board_list},
    numbers=left,
    basicstyle=\small\ttfamily,
    frame=single]
Board = [
  [x, '.', '.'],   % row 0
  ['.', '.', '.'], % row 1
  ['.', '.', o]    % row 2
]
% Symbols: '.' = free, x = P1 pawn,
%          o = P2 pawn, '#' = removed
\end{lstlisting}

% -------------------------------------------------------------------

\section{Report Structure}
\label{sec:structure}

This report is organised as follows:

\begin{itemize}
  \item \textbf{Chapter~\ref{chap:impl} --- Implementation:}
        Describes in detail all project modules: board representation,
        game rules, artificial intelligence algorithms, and the game
        interface.
  \item \textbf{Chapter~\ref{chap:eval} --- Experimental Evaluation:}
        Presents the results of the unit tests and performance
        benchmarks obtained under different configurations.
  \item \textbf{Appendix~\ref{app:code} --- Source Code:}
        Selected listings of the developed Prolog files.
\end{itemize}
